\scanpage{381}
\mspage{193}

\inserted{\unsure{germes de|germe de}} relèvements, pour \unsure{un|une} \inserted{germe d'}action extérieure sur $\pi$ déjà donné --- ce qui correspond au fait naïf que la courbe algébrique $U_{\deleted{\unsure{g,\nu}}}$ sur $\overline{\mathbb Q}$ définie par cette action extérieure \unsure{a} ``beaucoup de p\textsuperscript{ts}'' rat/$\overline{\mathbb Q}$. En fait, il suffirait, à la limite, de parler des \deleted{\illegible{1}} propriétés de ces relèvements plus complets en \inserted{un germe d'}une vraie action de \deleted{\illegible{1}} $\mathbb{\Gamma}_a$ sur $\pi$ (pas seulement extérieure) --- dont les actions extérieures \unsure{sont|vont} \unsure{à|se} déduire par composition avec
\[
M_a \longrightarrow N_a \simeq M_a/\pi.
\]

\deleted{\illegible{1}} On suppose donc donné un ss-groupe fermé
\[
\Gamma \subset \deleted{\unsure{\mathcal G}} M_a \subset \widehat{\widehat{\mathcal G}}(\pi)
\]
tel que
\[
\text{\textcircled{a}}\qquad
\Gamma \longrightarrow \mathbb{\Gamma}_a
 = M_a/\pi
\quad \text{\unsure{soit|est} injectif,}
\]
et \unsure{ait|a} comme image un ss-groupe ouvert de $\mathbb{\Gamma}_a$.

--- étant entendu que deux sous-groupes conjugués sous $\pi$ --- voire même, \inserted{parfois,} sous $M_a$ --- ne sont pas considérés comme essentiellement distincts. On considère, en $\overline{\mathrm m}$ temps que $\Gamma$, des sous-groupes

\scanpage{382}

ouverts $\Gamma'$. On veut sûrement, pour de tels sous-groupes ouverts
\[
\text{\textcircled{b}}\qquad \pi^{\Gamma'}=(1)
\]
Si on désigne par $\widetilde{\Gamma}$ l'image de $\Gamma$ dans $N_a$, de sorte qu'on a une extension
\begin{samepage}
(1)\quad
\diagram{\begin{tikzcd}[ampersand replacement=\&,column sep=2.1em,row sep=1.1em]
1 \arrow[r,dash] \& \pi \arrow[r,dash] \& \mathcal E \arrow[r] \& \widetilde{\Gamma} \arrow[r,dash] \& 1 \\
\& \& \& \Gamma \arrow[u,"\sim"'] \&
\end{tikzcd}}
\end{samepage}
on peut considérer $\Gamma$ comme une section de cette extension, $\Gamma'$ comme une section partielle. L'hyp. $\pi^{\Gamma'}=(1)$ assure la \uline{rigidité} de la catégorie des points de $U$ à valeurs dans $\overline{\mathbb Q}$, paradigmée par celle des germes de \unsure{scindages|scindage} de l'extension \textcircled{(1)'}.
\margin{\unsure{!|?}}

On peut aussi regarder les ensembles
\[
\unsure{\widehat{\mathcal G}_a^{+\Gamma'}}
=
\deleted{\unsure{N_a^{\Gamma'}|M_a^{\Gamma'}}}
\operatorname{Centr}_{M_a}(\Gamma')\cap\widehat{\mathcal G}_a^+
\quad\text{et}
\]
\[
\widehat{\mathcal T}_a^{+\Gamma'}
=\operatorname{Centr}_{N_a}(\widetilde{\Gamma}')\cap\widehat{\mathcal T}_a^+,
\]
\deleted{et} on a donc une suite exacte
\[
(2)\qquad 1\longrightarrow\pi^{\Gamma'}\longrightarrow
(\widehat{\mathcal G}_a^+)^{\Gamma'}\longrightarrow
(\widehat{\mathcal T}_a^+)^{\Gamma'}
\]
qui, compte tenu de l'hyp. $\pi^{\Gamma'}=(1)$,

\scanpage{383}
\mspage{194}

donne
\[
(3)\qquad (\widehat{\mathcal G}_a^+)^{\Gamma'}
\hookrightarrow (\widehat{\mathcal T}_a^+)^{\Gamma'}.
\]

Le deuxième membre de (3), par notre dictionnaire hypothétique, devrait être \inserted{\unsure{canoniquement}} \unsure{le|iso} groupe des \unsure{automs} de la courbe $U$ sur $\unsure{K'_a|K'}$ donc être un groupe fini, et $\overline{\mathrm m}$ \deleted{\unsure{par}} la limite inductive (pour $\Gamma'$ décroissant) --- qui est le groupe des \unsure{automs} de $U$ sur $\overline{\mathbb Q}$ --- est fini. Désignant par un exposant $\natural$ le germe de groupe correspondant, on veut donc
\[
\text{\textcircled{c}}\qquad
Z=(\widehat{\mathcal T}_a^+)^{(\Gamma^{\natural})}
\left(=\operatorname{Centr}_{\unsure{N_a|M_a}}
(\widetilde{\Gamma}^{\natural})\cap\widehat{\mathcal T}_a^+\right)
\]
est un groupe \uline{fini}

--- ce qui fait pendant à (b), et exprime que les groupes d'automorphismes des points de $M_{g,\nu\overline{\mathbb Q}}$ sur $\overline{\mathbb Q}$ sont finis.

En fait, on voudrait que $Z$ soit conjugué dans $\deleted{\unsure{\widehat{\mathcal T}}_a^+}\inserted{M_a}$ à un sous-groupe de \unsure{à} $\mathcal T^+$ \deleted{Centr} (si on suppose qu'on dispose d'une \inserted{discrétification} \deleted{\illegible{1}} $\pi_0$ de $\pi$, permettant de définir $\mathcal T$ --- sinon, \unsure{on peut} exprimer cette

\scanpage{384}

propriété en disant qu'il existe une \unsure{discrétification} \inserted{$\pi_0$} de $\pi$, compatible avec $a$, telle que $Z\subset\mathcal T(\pi_0)$\ldots)

\unsure{Considérons} maintenant un sous-groupe fini quelconque $G\subset\widehat{\mathcal T}_a^+$. [N.B. Si on prend \inserted{seulement} $G\subset M_a$, de sorte que l'image de $G$ dans $\mathbb{\Gamma}_a=M_a/\widehat{\mathcal T}_a^+$ est un ss-groupe fini, alors s'il \deleted{$\mathbb{\Gamma}_a$} est vrai que $\mathbb{\Gamma}_a\simeq\unsure{\mathcal G|\mathfrak G}$, \margin{$\operatorname{Gal}(\overline{\mathbb Q}/\mathbb Q)\longrightarrow$} cette image doit être d'ordre $1$ ou $2$, et dans le deuxième cas doit \margin{\unsure{?}} définir un $\tau\in\mathbb{\Gamma}_a$ correspondant à une \uline{\unsure{prédiscrétification}} (pas orientée) bien déterminée de $\pi$. Ce cas devrait être encore réutilisé dans la suite\ldots]. \deleted{Si on} \inserted{Supposons, \unsure{ou}} \unsure{admettons}, que $G$ est contenu dans $\mathcal T^+$, \margin{ce n'est sans doute \uline{pas} automatique (si on \unsure{ne} suppose \unsure{pas|N'} $N_a=\operatorname{Norm}_{\widehat{\widehat{\mathcal T}}}(\unsure{\widehat{\mathcal T}_a|\widehat{\mathcal T}_a^+})$)} pour une discrétification convenable dans la classe $a$, et que l'action extérieure ainsi obtenue de $G$ sur un $\pi_0$ discret de type $(g,\nu)$ \unsure{soit|sont} tjrs \uline{réalisable}. Elle est donc réalisable aussi \unsure{par|pour} une action de $G$

\scanpage{385}
\mspage{195}

sur une structure complexe, donc algébrique, ce qui signifie qu'on peut trouver un \deleted{\illegible{1}} \inserted{\unsure{germe|germes} de relèvement \uline{admissible}}
\[
\mathbb{\Gamma}_a\longrightarrow N_a
\]
qui centralise $G$. Considérons d'autre part
\[
(4)\qquad
\underset{\inserted{\widehat N_a}}{\operatorname{Centr}_{N_a}}(G)
\longrightarrow\mathbb{\Gamma}_a
\]
dont le noyau est $\operatorname{Centr}_{\widehat{\mathcal T}_a^+}(G)=\widehat{\mathcal T}_a^{+G}$.
Si on se \unsure{plaçait|plaçant} dans le contexte discret \deleted{\unsure{pour}} (avec une \inserted{discrétification} $\pi_0\subset\pi$ \unsure{invariante|invariants} par $G$) on aurait, par les conjectures standard topologiques, que $\unsure{\widehat{\mathcal T}_a^{+G}|\widehat{\widehat{\mathcal T}}_a^{+G}}$ est le groupe des autom., dans la catégorie isotopique, de l'action de $G$ sur une surface $U$ de type $(g,\nu)$, \unsure{décrite} par l'op. ext. donnée de $G$ sur $\pi_0$. Ce groupe n'a aucune raison d'être fini --- s'il l'était, l'\unsure{hom.} (4) serait à noyau fini, donc serait un \unsure{iso} \unsure{des|de} noyaux de groupes, et il y aurait un

\scanpage{386}

germe de relèvement unique $\mathbb{\Gamma}_a^{\natural}\longrightarrow M_a$ qui centraliserait $G$ --- ce qui signifierait qu'il y a \deleted{une seule} une seule façon de réaliser l'opération \inserted{topologique} de $G$ sur $U$ par une opération \deleted{algébrique} analytique complexe, donc algébrique --- or ce n'est surement pas le cas, l'\unsure{ens} des \unsure{pts} fixes de $G$ opérant sur l'espace de Teichmüller $\widetilde M_{g,\nu}$ n'est pas réduit à un point --- c'est \inserted{(dans le contexte topologique)} une multiplicité \unsure{schématique|arithmétique} qui peut être \deleted{\illegible{2}} de dim $\unsure{\mathbb R|\mathbb C}$ quelconque --- Elle est de dim $>0$ en tous cas, si le quotient \deleted{\illegible{1}} $U/G$ n'est pas de genre $0$.

À retenir en tous cas, comme propriété plausible :

(d) Pour $\forall$ ss-groupe fini \inserted{$G$} de $\widehat{\mathcal T}_a^+$ (ou du moins si $G\subset\mathcal T^+$, quand on dispose d'une \inserted{\unsure{discrét.}} $\pi_0\subset\pi$ \inserted{\unsure{dans $a$}}), regardant $\operatorname{Centr}_{\unsure{\widehat{\widehat{\mathcal T}}|\widehat{\mathcal T}}}(G)=\unsure{;}Z(G)$, $N_a\cap Z(G)\inserted{\longrightarrow\mathbb{\Gamma}_a}$ \deleted{est ouvert dans $N_a$} \inserted{a une image ouverte (et il y a $\unsure{\overline{\mathrm m}}$ des}

\scanpage{387}
\mspage{196}

germes de relèvements admissibles $\mathbb{\Gamma}_a\to Z(G)$)

Ceci implique une propriété non triviale pour
les $g\in\unsure{N_a|M_a}$ \illegible{1} en tant qu'éléments de $\widehat{\widehat{\mathcal T}}(\pi)$,
ou mieux de $\inserted{N'=}\operatorname{Norm}_{\widehat{\widehat{\mathcal T}}}(\widehat{\mathcal T})$ :
\deleted{$\exists n\in\mathbb N^{*}$ tel que $g^n$ \unsure{,|;}}
\margin{\unsure{savoir} : $\exists n\in\mathbb N^{*}$ tel que $g^n$ soit congru mod $\widehat{\mathcal T}^{+}_a$ à un élément de $Z(G)$\ldots}
Considérons en effet l'image \inserted{$\unsure{Z'_G|Z'(G)}$} de
$Z(G)\cap N'$ dans $\Gamma'=N'/\widehat{\mathcal T}$, \deleted{\unsure{alors est|alors soit}}
\inserted{où évidemment $\mathbb{\Gamma}_a\subset\Gamma'$, et \deleted{\unsure{la propriété}} l'image de $Z(G)\cap N_a$ dans $\mathbb{\Gamma}_a$ n'est autre que}
\deleted{\illegible{1} soit $\mathbb{\Gamma}_a$ celle de $\Gamma$}
$Z'_G\cap\mathbb{\Gamma}_a$. Dire que celle-ci est \inserted{\unsure{ouverte|ouvert} i.e. \uline{d'indice}} \unsure{finie|fini},
implique donc que $\forall g\in\mathbb{\Gamma}_a$, \deleted{\illegible{1}} $\exists n\in\mathbb N^{*}$
tel que $g^n\in Z'_G$. \deleted{On \unsure{À vrai dire}}

Notons que dans $\mathcal T(\pi_{g,\nu})$ il n'y a
qu'un nb fini de classes de conjugaison
de sous-groupes finis, donc si on
regarde leurs centralisateurs \inserted{$Z_G$} dans
$\widehat{\widehat{\mathcal T}}(\pi_{g,\nu})$, et leurs images dans
$\Gamma'_{g,\nu}=\operatorname{Norm}_{\widehat{\widehat{\mathcal T}}_{g,\nu}}(\widehat{\mathcal T}^{+}_{g,\nu})/\widehat{\mathcal T}^{+}_{g,\nu}$, on ne
trouve qu'un nombre fini de sous-groupes
de $\Gamma'_{g,\nu}$, soit $Z'_{g,\nu}$ leur intersection.
On voit donc que \unsure{ce|le} sous-groupe $\mathbb{\Gamma}_{g,\nu}$
de $\Gamma'_{g,\nu}$ est tel que $\mathbb{\Gamma}_{g,\nu}\cap Z'_{g,\nu}$ soit
ouvert dans $\mathbb{\Gamma}_{g,\nu}$

\scanpage{388}

Passons maintenant \unsure{à} des ss-groupes
finis de $\widehat{\widehat{\mathcal G}}(\pi)$ lui-$\overline{\mathrm m}$ --- ou du moins
de $M_a$ --- ou, ce qui revient presque au
même, de $\widehat{\mathcal G}^{+}_a$. On va, comme
tantôt, se borner à des $G\subset\mathcal G(\pi_0)$,
pour une \unsure{discrétification} convenable $\in a$,
car autrement il n'y aurait guère rien à
dire. On sait qu'alors $G$ est
cyclique --- \unsure{il|et} correspond à une action
de $G$ sur un $U$ topologique, \uline{avec}
\uline{pt fixe}. On peut la réaliser de
façon complexe, ce qu'on \unsure{exprimera|exprime} en
disant qu'il existe un relèvement
admissible $\mathbb{\Gamma}_a^{\natural}\longrightarrow\deleted{\unsure{\mathcal G}} M_a$, qui
centralise $G$. On sait d'ailleurs, si
$G\ne1$, que \deleted{$\unsure{\pi}$}
\[
(5)\qquad \pi^G=(1)
\]
(ou plutôt, on le sait dans le cas discret ---
on l'admet dans le contexte profini) ---

\scanpage{389}
\mspage{197}

ce qu'on peut \unsure{encore} interpréter comme une propriété de rigidité --- ainsi
\[
\operatorname{Centr}_{\unsure{\widehat{\widehat{\mathcal G}}}}(G)
=\widehat{\widehat{\mathcal G}}^{G}
\longrightarrow\widehat{\widehat{\mathcal T}}^{G}
\]
est \uline{injectif}, donc si on a un $\unsure{\Gamma''|\widetilde{\Gamma}}\subset\unsure{N_a|M_a}$ qui est dans l'image, alors il existe un relèvement $\unsure{\Gamma''|\widetilde{\Gamma}}\longrightarrow\widehat{\widehat{\mathcal G}}$ unique qui centralise $G$ i.e. tel que ce soit $\Gamma\longrightarrow\widehat{\widehat{\mathcal G}}^{G}$.

Mais en fait ce qu'on saura, pour un $\Gamma\subset M_a$ donné (\unsure{définissant} une courbe algébrique via son action ext. sur un $\pi_1$) c'est que $\Gamma\subset M_a^G$ (i.e. l'action extérieure commute à une action extérieure donnée de $G$ \inserted{i.e. $G$ opère sur la courbe alg.}) --- et on voudrait néanmoins en conclure que lorsqu'on a relevé $G\longrightarrow\widehat{\mathcal T}^{+}$ \unsure{en|ou} $G\longrightarrow\widehat{\mathcal G}^{+}$ (i.e quand on s'est donné un pt fixe de l'action de $G$ sur $U$): \unsure{alors|donc} l'action de $\Gamma^{\natural}$ se remonte \inserted{(i.e. $\Gamma^{\natural}\longrightarrow M_a^G$)} \unsure{automatiquement|automatique}, \inserted{en $\Gamma^{\natural}\longrightarrow M_a$} de façon \unsure{à} \unsure{commuter|commutante|commutant},\ldots . Est-ce là une propriété \unsure{du choix des|des choix des} $M_a$, ou de celui du relèvement $\Gamma^{\natural}\unsure{\longrightarrow|\hookrightarrow}N_a$, ou de $G$??

\scanpage{390}

Dans le cadre discret, sauf erreur, si on a une action fidèle \inserted{\unsure{directe|discrète},} de $G$ fini sur $\pi_0$, alors \deleted{Centr}
\[
\mathcal G(\pi_0)^G\longrightarrow\mathcal T(\pi_0)^G
\]
(qui est injectif, par $\pi_0^G=1$) est aussi surjectif. Non, il y a erreur --- ça \unsure{signifierait} tel quel que si $G$ opère fid. sur une surface top. $U$ de type $g,\nu$, avec un pt fixe $P$, alors les autom. qui commutent à $G$ \uline{fixent} $P$ (au lieu de permuter seulement les pts fixes entre eux\ldots). Donc il ne faut pas s'attendre à ce que $\forall$ \unsure{élémt} de $\unsure{N_a^G|M_a^G}$, ni $\overline{\mathrm m}$ de $\widehat{\mathcal T}_a^G$, se remonte à $M_a^G$ --- mais plutôt ceci: pour un $G\unsure{\in|\subset}\mathcal T_a(\pi_0)$ \unsure{ss-groupe} fini donné, les classes de $\pi$-conj. de remontages à $\widehat{\mathcal G}$, \deleted{\unsure{à} $\mathcal G$} qui s'interprètent comme des pts fixes d'une action de $G$ sur quelque $U$, forment un ens. fini, sur lequel $\widehat{\widehat{\mathcal T}}(\pi)^G$ opère de façon naturelle, et \unsure{les|le} \unsure{stabilisateurs|stabilisateur} d'un pt $P$ i.e d'une classe de conj. de \unsure{rel\textsuperscript{ts}|relts}
